\section{Introduction}\label{sec:intro}

Supreme Court votes are the most consequential binary predictions in
American public life, and they are made under conditions that should
humiliate any forecaster: the inputs are thousands of pages of legal
argument, the decision rule is nine human beings, and the outcome
distribution is close enough to a coin flip that luck is difficult to
rule out. Yet a long line of work has shown that the task is not
entirely random. Statistical models built on case attributes have
predicted affirm/reverse outcomes at rates far above chance for two
decades~\citep{ruger2004,katz2017}, and the dominant theoretical
account of the Court---the attitudinal model~\citep{segal2002} ---
holds that ideology, not legal doctrine, explains most of the variance
in votes.

This paper reports the design and the first results of \emph{Legally
Subjective}, an independently conducted, zero-budget experiment that
asks a narrower and, we argue, more honest version of the prediction
question. Rather than asking how accurately a court's output can be
retrofitted by a flexible learner with centuries of data, we ask two
questions that can both fail loudly. First, is each justice's vote
predictable from the public case file alone, on a fixed, modern window,
against a pre-registered bar? Second---and this is the experiment's
live wire---if a language model is fine-tuned on one justice's own
past opinions, does it predict that justice's \emph{future} votes
better than the same model without the persona? Both answers are
publishable. If the persona helps, then judicial subjectivity leaves a
measurable textual footprint, and a justice's written voice carries
predictive information about their future decisions. If it does not,
then whatever is predictable about a justice is already captured by
their ideological direction, and the remainder is either private
deliberation or noise.

Three design commitments distinguish the project from prior work.
First, everything is pre-registered before the models see the
\emph{sealed} holdout: the baselines, the split, the metrics, and the
fifty-case test set are frozen in the repository, with the selection
seed and the SHA-256 of the sealed list published in advance. There is
no second holdout and no post hoc tuning. Second, the experiment is
built to be audited by a stranger: the corpus is assembled from public
sources with every raw file hash-chained, and every number in this
paper regenerates from scripts in the repository. Third, the budget is
zero by construction---public data, free API tiers, and free-tier
compute---which is simultaneously a constraint, a proof of
accessibility for non-institutional researchers, and a defense against
the critique that results were bought rather than earned.

The contributions of this working paper are: (i) a frozen,
provenance-complete corpus of 569 argued Supreme Court cases spanning
OT2015 through OT2023, fusing CourtListener metadata with Supreme Court
Database (SCDB) vote coding, with 98.6\% oral-argument coverage for
future multimodal conditions; (ii) pre-registered statistical
baselines, of which the strongest---per-justice modal ideological
direction---sets a vote-level bar of 63.7\% on the held-out window;
(iii) the first trained challenger, a family of structured-feature
models whose \emph{null result} is itself informative: nothing in the
pre-decision metadata adds stable signal beyond who the justice is
(Section~\ref{sec:results}); and (iv) the sealed, hash-frozen test
protocol under which the remaining conditions (zero-shot, persona
fine-tuning, retrieval) will eventually be judged
(Section~\ref{sec:protocol}).

We emphasize what this paper does not claim. It does not claim that
Supreme Court votes are unpredictable; two decades of literature and
our own baselines say otherwise. It does not report any language-model
results; the opinion-text corpus required for those conditions is
still being collected under free-tier API quotas. And it does not
treat the 63.7\% bar as trivial: as Section~\ref{sec:results} shows,
three families of supervised models failed to clear it on the same
rows, which is precisely why the bar is interesting.
